%==============================================================================
\section{System Overview}
%==============================================================================

\subsection{What the system does}
A Holter monitor is a portable device that records a patient's ECG continuously
while they go about normal activity. This project implements a modern, networked
variant: instead of storing the signal on the device for later review, the
wearable node \emph{streams} the heart waveform over Wi-Fi so that a clinician
can watch it live on a remote dashboard from anywhere on the network.

The patient wears a small enclosure containing an ESP32 microcontroller and an
\textbf{AD8232} single-lead ECG analog front-end. Three electrodes on the body
feed the heart's electrical activity into the AD8232, which conditions it into a
clean analog voltage. The ESP32 samples that voltage with its on-chip
analog-to-digital converter (ADC) and pushes batches of samples to a server. The
server filters the raw signal and broadcasts the result to every connected
browser, where it is drawn as a scrolling waveform.

\subsection{High-level architecture}
The system is split into three tiers, summarised in
Figure~\ref{fig:architecture}.

\begin{description}[leftmargin=2.2cm,style=nextline]
  \item[Acquisition] The ESP32 + AD8232 wearable node. Reads the analog ECG on
        \texttt{GPIO36} at $\approx$\,250\,Hz and transmits 250-sample JSON
        batches (one batch $\approx$ one second) over a Wi-Fi WebSocket.
  \item[Processing] A Django backend running on the Channels/Daphne ASGI stack.
        A WebSocket consumer receives each batch, runs an eight-stage DSP filter,
        and fans the cleaned signal out to all dashboard clients.
  \item[Visualisation] A single-page browser dashboard using Chart.js. It keeps a
        rolling 1000-sample buffer and redraws the live trace as data arrives.
\end{description}

\begin{figure}[h]
  \centering
  \begin{tikzpicture}[node distance=16mm and 20mm]
    % nodes
    \node[hw] (pat) {Patient\\\scriptsize 3 electrodes};
    \node[hw, right=of pat] (afe) {AD8232\\\scriptsize analog front-end};
    \node[hw, right=of afe] (esp) {ESP32\\\scriptsize ADC + Wi-Fi};
    \node[bk, below=18mm of esp] (srv) {Django + Channels\\\scriptsize ECGConsumer};
    \node[bk, left=of srv] (dsp) {DSP filter\\\scriptsize filter\_ecg\_signal()};
    \node[fe, left=of dsp] (web) {Browser\\\scriptsize Chart.js dashboard};

    % flows
    \draw[flow] (pat) -- (afe) node[midway,above,lbl]{mV};
    \draw[flow] (afe) -- (esp) node[midway,above,lbl]{0--3.3\,V};
    \draw[flow] (esp) -- (srv)
       node[midway,right,lbl,align=left]{Wi-Fi WebSocket\\\texttt{\{"ecg":[...]\}}};
    \draw[flow] (srv) -- (dsp);
    \draw[flow] (dsp) -- (web)
       node[midway,above,lbl,align=center]{WebSocket\\\texttt{\{"filtered":[...]\}}};

    % grouping backgrounds
    \begin{scope}[on background layer]
      \node[draw=hwcol!80!black, dashed, rounded corners, fit=(pat)(afe)(esp),
            inner sep=4mm, label={[font=\scriptsize\bfseries]above:Wearable node}] {};
      \node[draw=bkcol!70!black, dashed, rounded corners, fit=(dsp)(srv),
            inner sep=4mm, label={[font=\scriptsize\bfseries]below:Server (PC, port 8000)}] {};
    \end{scope}
  \end{tikzpicture}
  \caption{End-to-end architecture. The wearable node digitises and transmits raw
  ECG; the server filters it and broadcasts the cleaned waveform to the dashboard.}
  \label{fig:architecture}
\end{figure}

\subsection{Technology stack}
\begin{table}[h]
  \centering
  \small
  \begin{tabular}{@{}lll@{}}
    \toprule
    \textbf{Tier} & \textbf{Technology} & \textbf{Role} \\
    \midrule
    Acquisition  & ESP32 (Arduino C++)        & Sampling and Wi-Fi transport \\
                 & AD8232                     & Single-lead ECG analog front-end \\
                 & ArduinoWebsockets, ArduinoJson & WebSocket client + JSON encoding \\
    Processing   & Python 3, Django 5.1.7     & Web application framework \\
                 & Django Channels + Daphne   & ASGI / WebSocket handling \\
                 & NumPy, SciPy, PyWavelets   & ECG signal-processing pipeline \\
    Visualisation& HTML5 + JavaScript         & Single-page dashboard \\
                 & Chart.js 4.4.1             & Real-time line chart \\
    \bottomrule
  \end{tabular}
  \caption{Components used at each tier of the system.}
  \label{tab:stack}
\end{table}

\subsection{Repository layout}
The relevant source lives in two top-level folders:
\begin{lstlisting}[style=base,numbers=none]
HolterMonitor/
|-- arduino/
|   `-- websocket_ecg_send_data/
|       `-- websocket_ecg_send_data.ino   <- ESP32 firmware
`-- holter_monitor/                       <- Django project
    |-- manage.py
    |-- send_fake_data.py                 <- hardware-free ECG simulator
    |-- holter_monitor/                   <- project package
    |   |-- settings.py  asgi.py  urls.py
    `-- ecg/                              <- the ECG app
        |-- consumers.py   <- WebSocket consumer (ECGConsumer)
        |-- routing.py     <- ws/ecg/ route
        |-- filters.py     <- 8-stage DSP pipeline
        |-- views.py       <- serves the dashboard page
        `-- templates/ecg.html            <- Chart.js dashboard
\end{lstlisting}
